% A141: remove the actual boundary inertia from A139.
\subsection{A proper-curve detector of order \(512\) and class six}

\begin{statement}{A.141}{Proved}
{The same completed pair is detected after removing boundary inertia}
Retain the actual degree-eight curve \(D^\circ\), its proper model
\(\overline D\), the geometric groups \(\Pi,H_D,K_D\), and the completed
pair of A.139.  Let
\[
 \begin{gathered}
 \mathcal I=\ker\bigl(H_D\longrightarrow
                         G_D:=\pi_1(\overline D_{\overline\Q})\bigr),\\
 \overline K_D=
   \ker\bigl(H_D\longrightarrow G_D^{\rm ab}/2\bigr)
       =K_D\mathcal I .
 \end{gathered}
\tag{A.141.1}
\]
The product on the right is a closed subgroup.
It is normal in the full arithmetic group \(E_U\).
There is an actual geometric quotient
\[
 \overline Q_D=\Pi/\overline K_D
       \simeq
       \bigl(\F_2[\tau]/(\tau^6),+\bigr)\rtimes C_8,
       \qquad |\overline Q_D|=512,
\tag{A.141.2}
\]
where the cyclic generator acts by multiplication by \(1+\tau\).
Its nilpotency class is exactly six.

There is no section of \(E_U/\overline K_D\to G_\Q\) with the
combined mod-eight Kummer class (A.139.5), with the ordinary local
images \(x_p\) at all finite primes, and with an ordinary or cuspidal
real image.  In particular the same pair is excluded through the
actual maps
\[
    E_U/\Pi''\longrightarrow E_U/\overline K_D,\qquad
    E_U/\Gamma_7\Pi\longrightarrow E_U/\overline K_D.
\tag{A.141.3}
\]
The locally matching metabelian and class-six conclusions refer to
this explicit pair; they do not impose such a conclusion on all
infinite-support pairs.
\end{statement}

\begin{proof}
\emph{The proper comparison and the geometric degree correction.}
The proper curve \(\overline D\) has genus three and the rational
smooth point
\[
                   P_0:\quad T=W=0.
\tag{A.141.4}
\]
Indeed the derivative in \(T\) of \(F(T)-aW^4\) is \(1/2\) there.
The natural maps give the commutative arithmetic diagram
\[
 \begin{array}{ccccccccc}
 1&\longrightarrow&H_D&\longrightarrow&E_D&\longrightarrow&G_\Q&\longrightarrow&1\\
 &&\downarrow&&\downarrow&&\Vert\\
 1&\longrightarrow&G_D&\longrightarrow&\overline E_D&\longrightarrow&G_\Q&\longrightarrow&1,
 \end{array}
 \qquad
 \overline E_D=\pi_1(\overline D),\quad E_D=\pi_1(D^\circ).
\tag{A.141.5}
\]
Both vertical maps are onto, with kernel \(\mathcal I\).
The finite-cover description proves this: a finite etale cover
of \(\overline D\) is exactly a finite etale cover of \(D^\circ\)
whose four boundary inertia groups act trivially.

We use the proper comparison
\[
 H^2(\overline E_D,\mu_2)
                  \xrightarrow{\ \sim\ }
 H^2_{\rm et}(\overline D,\mu_2).
\tag{A.141.6}
\]
It is the actual natural group/etale map; it commutes with the
geometric restriction and with all point pullbacks.
An opened primary statement supplying it is
\cite[Theorem 1.1, Lemma 3.3 and Theorem 3.7]{LegLevrat2024}.
The geometric mechanism in this case is short.
Every connected finite geometric cover \(Y\) of \(\overline D\)
still has positive genus.  Its degree-one classes with coefficients
\(\F_2\) die after pullback to their finite quadratic torsors.
Its degree-two group is \(\F_2\), and a connected double cover
kills it because top-degree pullback multiplies by two.
Such a double cover exists since \(H^1(Y,\F_2)\) has positive
dimension.  Higher geometric cohomology vanishes
\cite[Propositions 14.1 and 14.2]{AgentMilneLEC}.
Thus the direct limit of the positive-degree cohomology over
finite geometric covers is zero.  The covering
Cartan--Leray comparison gives the geometric group/etale
isomorphism, and the two arithmetic Hochschild--Serre spectral
sequences give (A.141.6).  This also explains the precise maps in
the cited comparison.

Choose a lift
\(\widetilde\alpha\in H^2_{\rm et}(\overline D,\mu_2)\)
of the proper class \(\alpha_D\) from A.139.
On the geometric curve, the degree map identifies
\[
 H^2_{\rm et}(\overline D_{\overline\Q},\mu_2)=\F_2,\qquad
                   c_1(\mathcal O(P_0))\longmapsto1.
\tag{A.141.7}
\]
If \(\epsilon\in\F_2\) is the geometric restriction of
\(\widetilde\alpha\), replace it by
\[
              \widetilde\alpha_0=
                  \widetilde\alpha-\epsilon c_1(\mathcal O(P_0)).
\tag{A.141.8}
\]
The Kummer sequence shows that this has the same Brauer image
\(\alpha_D\), while its geometric restriction is zero.
It also has the same pullback at every rational point over any
field extension: a line bundle on a field is trivial, so its
first Chern class pulls back to zero.
The correction is made on the proper curve, and therefore has
the same property for a tangential section after its map to a
proper boundary point.

\emph{Exact inflation through proper mod-two homology.}
Let \(e_0\in H^2(\overline E_D,\F_2)\) correspond to
\(\widetilde\alpha_0\) by (A.141.6), and represent it by a
continuous central extension
\[
 1\longrightarrow\F_2\longrightarrow
      \widetilde{\overline E}_D\longrightarrow
      \overline E_D\longrightarrow1.
\tag{A.141.9}
\]
Its restriction class on \(G_D\) is zero by (A.141.8).
The definition of continuous group cohomology therefore supplies
a continuous splitting \(j:G_D\to\widetilde{\overline E}_D\).
For \(z\in\widetilde{\overline E}_D\), the two lifts of conjugation
on \(G_D\),
\[
       g\longmapsto zj(g)z^{-1},\qquad
       g\longmapsto j(\overline z g\overline z^{-1}),
\tag{A.141.10}
\]
differ by a continuous character \(G_D\to\F_2\).
It kills
\[
                     L_D=\overline{G_D^2[G_D,G_D]}.
\tag{A.141.11}
\]
Hence \(j(L_D)\) is closed and normal in the extension, and
intersects its central kernel trivially.
Exactly as in (A.137.9), the explicit map
\[
 \begin{gathered}
 \widetilde{\overline E}_D
   \longrightarrow
   \overline E_D\mathop{\times}_{\overline E_D/L_D}
               \bigl(\widetilde{\overline E}_D/j(L_D)\bigr),
       \qquad z\longmapsto(\overline z,[z])
 \end{gathered}
\tag{A.141.12}
\]
is an isomorphism of profinite groups.
For surjectivity, multiply a representative of the second
coordinate by the required \(j(l)\), \(l\in L_D\), to obtain
the first coordinate.  Injectivity follows from
\(j(L_D)\cap\F_2=1\); compactness gives the topology.
This proves exact inflation of \(e_0\) from
\(H^2(\overline E_D/L_D,\F_2)\).

The inverse image of \(L_D\) in \(H_D\) is \(K_D\mathcal I\),
because squares and commutators map onto squares and commutators
under the surjection in (A.141.5).
The product is closed as the product of two compact normal
subgroups.  Consequently the actual diagram identifies
\[
        E_D/\overline K_D
            \xrightarrow{\ \sim\ }\overline E_D/L_D .
\tag{A.141.13}
\]
This gives the desired cohomology class on the indicated cover
subgroup of \(E_U/\overline K_D\).
No class on the whole larger arithmetic quotient is claimed or
needed.

\emph{Boundary inertia, arithmetic normality and the finite group.}
Use the actual meridians \(x,y\) and the Nielsen coordinates
\(u=xy^{-2},v=y\) of A.139.
The four boundary inertia groups of \(D^\circ\) have generators
conjugate to
\[
                   x^4,\quad yx^4y^{-1},\quad
                   y^8,\quad (xy)^8.
\tag{A.141.14}
\]
At infinity the inverse is an equally valid generator.
The ramification indices are respectively \(4,4,8,8\),
and these are all four boundary points.
Their closed normal closure in \(H_D\) is exactly
\(\mathcal I\), by (A.141.5).
It is also stable under \(\Pi\): deck transformations permute
the boundary points.
The arithmetic action sends a geometric inertia generator
over each rational base cusp to a conjugate of its cyclotomic
unit power; its indicated powers generate the same closed
subgroups.  It therefore preserves \(\mathcal I\).
Together with the arithmetic normality of \(K_D\), this proves
that \(\overline K_D\) is normal in \(E_U\).

Write \(z=v^8\) and regard \(u_0,\ldots,u_7,z\) as the nine
independent order-two generators of \(H_D/K_D\).
The actual images of (A.141.14) are
\[
 \begin{array}{c|c}
 x^4&u_0+u_2+u_4+u_6+z\\
 yx^4y^{-1}&u_1+u_3+u_5+u_7+z\\
 y^8&z\\
 (xy)^8&u_0+u_1+\cdots+u_7+z .
 \end{array}
\tag{A.141.15}
\]
These generate the three-dimensional subspace spanned by
\(z\), the even lamp sum, and the odd lamp sum.
It is stable under \(v\), and hence is their full normal closure.
In the lamp module \(\F_2[C_8]=\F_2[\tau]/(\tau^8)\),
\(\tau=v-1\), the even sum is
\[
             1+v^2+v^4+v^6=(v-1)^6=\tau^6,
\tag{A.141.16}
\]
and the odd sum is \(v\tau^6\).
Killing these sums gives \(\F_2[\tau]/(\tau^6)\);
killing \(z\) replaces \(C_{16}\) by \(C_8\).
This proves the isomorphism (A.141.2), with the actual maps
\[
                           x\longmapsto u_0v^2,\qquad
                           y\longmapsto v .
\tag{A.141.17}
\]
The order is \(2^6\cdot8=512\).
The recurrence of A.139 becomes
\[
 \Gamma_k\overline Q_D=\tau^{k-1}
                  \F_2[\tau]/(\tau^6)\quad(2\le k\le6),
                 \qquad \Gamma_7\overline Q_D=1 .
\tag{A.141.18}
\]
In particular the weight-six commutator obtained from \([x,y]\)
by four more commutators with \(y\) is
\(\tau^5u_0=u_0+u_1+u_4+u_5\ne0\).
This proves exact class six.
The group remains metabelian, and
\(\Pi''\subset K_D\subset\overline K_D\).
These are precisely the kernels needed for (A.141.3).

The combined character \(2\kappa(t)+\kappa(1-t)\) has values
\(8,8,8,24\), up to sign, on the four boundary generators.
It is consequently zero modulo eight on \(\mathcal I\),
as well as on \(K_D\).  Thus the combined class in (A.139.5)
is still defined on \(E_U/\overline K_D\), and its zero
root-torsor class still gives a fixed sheet of the actual
degree-eight cover.

\emph{The local contradiction through this quotient.}
A section with the stated data therefore lifts, after geometric
conjugation, to \(E_D/\overline K_D\).
Pull back the class from (A.141.12)--(A.141.13).
At each finite prime, the ordinary local comparison supplies a
fixed sheet in the actual cover above \(x_p\), hence a
\(\Q_p\)-point of \(D^\circ\).  The class evaluates there as
\(\alpha_D\), since the correction (A.141.8) has zero point
pullback.  A.139 considered every possible such sheet and proved
invariant \(1/2\) at two and zero at every other finite prime.
The real ordinary image evaluates at a real point.
For a real cuspidal image, the fixed sheet gives a real inertia
orbit, hence a real boundary point of \(\overline D\).
The proper cohomology map kills the boundary inertia, so this
pullback is again the value at that proper point.
It is zero because \(d>0\).

We have obtained a class in \(\operatorname{Br}(\Q)[2]\)
whose local invariants sum to \(1/2\), contradicting global
reciprocity.  This proves the finite-quotient exclusion and all
its projected consequences.
\end{proof}

The standard-library certificate
\path{code/harmonic_proper_detector_certificate.py},
its output and
\path{code/harmonic_proper_detector_data.json}
verify all four literal boundary words, the order-eight boundary
kernel, every one of the \(4096\) original elements and its image
in the \(512\)-element quotient, and \(12288\) generator-product
images.  They check the lower-central dimensions
\(5,4,3,2,1,0\), the displayed nonzero weight-six word, and all
\(262144\) products for the retained mod-eight character.
